%!TEX root = ../../thesis.tex

\subsection{libxml2}\label{sec:target-libxml2}

% GROUND TRUTH (~/fuzzyemu/sysroots/resources/cases/libxml2/main.c;
% figures/generated/tab-campaign-inventory.tex):
%   - guest: persistent-mode harness reading a fixed-size record from stdin, calling
%     xmlReadMemory on exactly that many bytes, freeing the tree and looping.
%     xmlInitParser is called before the snapshot point so the one-time parser setup
%     is not paid inside the first measured parse.
%   - parse options: XML_PARSE_RECOVER | XML_PARSE_NOERROR | XML_PARSE_NOWARNING.
%   - campaign LIBXML2-400: one arm, 4 seeds, 400 iterations, T = 2,
%     xi = 0.03906, median 5281 blocks (5085--5419), ~42 h per run.
%   - best input: `<:>/..<P.&.<:....4<b...&H:..<...`, 4382 blocks (32 bytes).
%   - structural byte rate series: `xml_bytes`.

% TODO: The largest target in the thesis by an order of magnitude --- open with that,
% because scale is the reason it is here. It shares the corpus MDP, the harness shape,
% the reward and the instrumentation with \cref{sec:target-cjson}; describe only what
% differs, and \cref the shared machinery rather than restating it.
%
% What differs, and why each difference matters:
%   - the guest is a real, widely deployed XML stack rather than a single-file parser.
%     ~5300 blocks are reached against cJSON's ~900, which is what makes this the
%     target where the hashed bucket array of \cref{sec:platform-coverage} is actually
%     stressed. Report the collision rate at this block count; if the map saturates,
%     that is a headline limitation and not a footnote (\cref{sec:discussion-threats}),
%   - XML_PARSE_RECOVER is a deliberate harness choice with a coverage consequence: it
%     keeps the parser walking past a well-formedness error and returns the partial
%     tree, so a mutation that breaks well-formedness still yields coverage distinct
%     from its neighbours instead of collapsing every malformed record onto one
%     early-abort path. State this as an experimental decision that shapes the reward
%     landscape --- it makes the target more gradient-rich than a strict parser would
%     be, which cuts both ways for the conclusions,
%   - the record is longer than cJSON's, so the place phase ranges over more byte
%     positions while |A| stays 256.
%
% Its role: transfer. The configuration is not tuned here; it is inherited from the
% cJSON ablation and run unchanged. That is the substitute for a train/test split over
% environments (\cref{sec:methodology-hyperparameters}), so the section must be
% explicit that nothing about this target informed the configuration.
%
% What the target tests that cJSON cannot: whether coverage-driven mutation planning
% survives a target where the reachable block count is far larger than the agent can
% enumerate within an episode, and whether XML's nesting structure is discoverable from
% raw coverage the way JSON's bracket nesting is. Compare the best inputs of the two
% targets directly (\cref{tab:best-inputs}) --- both converge on deep structure, and
% that is the most legible evidence in the thesis that the agent found the grammar.
%
% OPEN ITEMS:
%   - confirm the record length actually used (the best input is 32 bytes, while the
%     sysroots harness in the working tree is compiled at INPUT_LENGTH 16). The
%     campaign revisions are not in the local checkout; read the harness of
%     \texttt{18a8f75}/\texttt{fc8884b} before writing a number.
%   - the four seeds ran on TWO revisions (\texttt{18a8f75} and \texttt{fc8884b}), so
%     they are not repeated samples of one configuration. Say so here, not only in the
%     inventory footnote, and say what changed between them.
%   - there is no random control arm for this target. \cref{fig:coverage-executions}
%     is therefore read against the structural-rate chance level of
%     \cref{fig:structural-rate} instead; state that substitution and its weakness.
